\section{Доказательства свойств для множества Мандельброта}

\subsection{Множество Мандельброта переходит само в себя при сопряжении}

Пусть \(c \in M\) ,\,последовательность \(z_{n+1} = z_n^2 + c, z_0 = 0\) ограничена.

Рассмотрим последовательность для \(\bar{c}\):
\[
w_0 = 0, \quad w_{n+1} = w_n^2 + \bar{c}
\]

Докажем индукцией, что \(w_n = \overline{z_n}\) для всех \(n\):

\(n=0\): \(w_0 = 0 = \overline{0} = \overline{z_0}\)\\
\(w_n = \overline{z_n}\). Тогда:
  \[
  w_{n+1} = w_n^2 + \bar{c} = (\overline{z_n})^2 + \bar{c} = \overline{z_n^2 + c} = \overline{z_{n+1}}
  \]

Так как \(|w_n| = |\overline{z_n}| = |z_n|\), то ограниченность \(\{z_n\}\) влечёт ограниченность \(\{w_n\}\). 

Следовательно, \(\bar{c} \in M\). 

\subsection{Если \(|c|>2\), то \(c\) не принадлежит множеству Мандельброта}

Рассмотрим последовательность Мандельброта:
\[
z_0 = 0, \quad z_{n+1} = z_n^2 + c
\]
Предположим, что \(|c| > 2\).

\[
|z_1| = |{z_0}^2 + c | = | 0 + c | = |c| > 2
\]

Пусть для некоторого \(n \geq 1\) верно:
\[
|z_n| \geq |c| \quad \text{и} \quad |z_n| > 2
\]

\[
|z_{n+1}| = |z_n^2 + c| \geq |z_n|^2 - |c| \quad
\]
По предположению \(|z_n| \geq |c|\), следовательно:
\[
|z_{n+1}| \geq |z_n|^2 - |z_n| = |z_n|(|z_n| - 1)
\]
Так как \(|z_n| > 2\), то \(|z_n| - 1 > 1\), поэтому:
\[
|z_{n+1}| > |z_n|
\]
Кроме того, так как \(|z_n| \geq |c|\) и \(|z_n| > 2\), то:
\[
|z_{n+1}| \geq |z_n|(|z_n| - 1) > |z_n| \geq |c|
\]

Мы показали, что для всех \(n \geq 1\):
\begin{itemize}
    \item \(|z_n| \geq |c| > 2\)
    \item \(|z_{n+1}| > |z_n|\)
\end{itemize}

Следовательно, \(|z_n| \to \infty\), и \(c\) не принадлежит множеству Мандельброта при \(|c| > 2\). 

